\chapter{\abstractname}

Contract Lifecycle Management (CLM) platforms increasingly rely on large language
models (LLMs) to automate clause extraction from commercial contracts, and the research
and vendor communities have advocated multi-agent architectures as a superior
alternative to single-prompt deployments. Despite this advocacy, no published work has
rigorously compared multi-agent architectural patterns against single-agent baselines
across multiple frontier-model families on a standardised contract-analysis benchmark
under realistic class distributions.

This thesis presents a controlled, statistically tested comparison of five
architectural configurations --- zero-shot (B1), chain-of-thought (B4), a
classify-then-extract pipeline (M0), an LLM-routed specialist system (M1), and a
consolidated specialist ablation (M6) --- evaluated across eleven frontier models from
the Anthropic, Google, and OpenAI families on 81 long contracts from the CUAD benchmark.
Evaluation uses a natural $30\%$ positive / $70\%$ negative class distribution, $F_2$
and $F_1$ as primary metrics, and paired McNemar and bootstrap statistical tests with
Benjamini--Hochberg correction.

The central finding is that no multi-agent configuration consistently outperforms the
strongest single-agent baseline on $F_2$: pooled mean $F_2$ across all eleven models
falls within a 1.1-point band ($\mathrm{B1} = 0.801$, $\mathrm{M0} = 0.794$,
$\mathrm{M1} = 0.790$). Architectural ranking varies systematically by model family.
Two mechanistic findings explain the null result. First, the classify-then-extract
pipeline (M0) suffers a \emph{classification-gate rejection cascade} on Gemini-family
models, with wrong-rejection rates of $31$--$41\%$ versus $1.5$--$2.7\%$ for the
single-agent baseline, accounting for $70$--$98\%$ of M0's false negatives. Second, the
LLM router in M1 achieves near-perfect routing accuracy ($100\%$ on ten of eleven
models), collapsing to a trivially solvable category-name lookup; M1 is statistically
indistinguishable from the consolidated single-agent ablation M6 on $F_2$
($\Delta F_2 = 0.000$), implying that the architectural decomposition contributes
nothing beyond the prompt content. Multi-agent configurations additionally carry a
$24\%$--$70\%$ provider-cost premium with no aggregate accuracy benefit, placing the
zero-shot single-agent baseline on the Pareto frontier.

These results delineate the conditions under which multi-agent designs help, are
inert, or actively hurt in enterprise CLM, and motivate a practitioner-oriented
six-gap analysis connecting benchmark performance to production deployment reality.
